\section{Physical Limitations and Alternative Topologies}
\label{sec:limitations}

While the splicing mechanism of \textsc{Nexus} eliminates the online prefill wall under standard workloads, it encounters physical and mathematical boundaries when applied to deep contexts and multi-tool attention matrices. In this section, we analyze these limits and evaluate alternative caching topologies.

\subsection{Zero-Copy Memory Bandwidth Constraint}
\label{sec:zero-copy-bandwidth}

Rotary Position Embeddings (RoPE) represent token positions as rotation angles in a multi-dimensional complex plane. Splicing a pre-compiled KV cache block (compiled starting at position $m_0=0$) into a live context at depth $P$ alters the positional coordinates. If $P \le 256$, the resulting RoPE phase drift is minor, and recomputing a small suffix fraction (e.g., $5\%$) is sufficient to restore the attention distribution, yielding a low KL divergence of $0.0076\,\text{nats}$ and a perfect $1.0$ Top-1 logit agreement.

At deeper context sizes ($P = 1024$), however, standard splicing strategies degrade because the absolute phase drift accumulates nonlinearly: the phase error $\Delta \theta_i = (n_{\text{past}} - m_0)\theta_i$ shifts key representations significantly. While this drift is mathematically correctable by applying a relative rotation matrix to the Key tensors (for instance, via the \texttt{llama\_kv\_cache\_seq\_shift} kernel), executing such an $\mathcal{O}(P \cdot d_{head})$ read-modify-write kernel over the KV cache would trigger VRAM copies. This execution completely invalidates the core $\mathcal{O}(1)$ zero-copy UMA \texttt{MAP\_SHARED} host-pointer paradigm of \textsc{Nexus}. 

Thus, this constraint is not an absolute physical barrier, but a strict memory-bandwidth tradeoff: the overhead of executing online relative rotations over deep caches outweighs the savings of avoiding the prefill computation, forcing the system to fall back to a full Python text prefill (Path B) for deeper context sequences.

\subsection{The Causal Isolation Constraint}
\label{sec:causal-isolation}

To support queries requiring the coordination of multiple independent tools, we evaluate a multi-tool cache topology. In this setup, we compile the KV caches of multiple tools separately and attempt to splice them concurrently into the active context. 

To prevent cross-tool token contamination, we implement \textbf{Multi-Tool Blockmasking ($N1m$)}. The self-attention matrix is partitioned into block-diagonal segments, where each tool block is causally blinded from the token coordinates of adjacent tools.

\Cref{tab:multi_tool_paradox} (derived from our multi-tool evaluation suite) summarizes the performance of this topology:
\begin{table}[h]
\centering
\caption{Multi-tool topology performance comparison.}
\label{tab:multi_tool_paradox}
\small
\begin{tabular}{lrrr}
\toprule
Topology & Tensor KL & E2E Accuracy & Status \\
\midrule
Single Tool Splice ($N1$) & 0.0076 & 0.91 & Sound \\
Multi-Tool Blockmask ($N1m$) & 0.0000 & 0.00 & Blinded \\
\bottomrule
\end{tabular}
\end{table}

The empirical evaluation exposes the \textbf{Causal Isolation Constraint}. Under multi-tool blockmasking ($N1m$), the attention calculations yield a perfect tensor KL divergence of $0.0000\,\text{nats}$ relative to isolated reference runs. However, this causal isolation completely blinds the tools from each other during the prefill phase, preventing the model from building a calibrated joint latent distribution. This blinding ruins the cross-tool routing probability, and the end-to-end routing accuracy drops to exactly $0.0\%$.

This result demonstrates a critical systems insight: minimizing tensor-level KL divergence is an insufficient proxy for task-level accuracy. Causal interaction across tool schemas is not a source of noise to be masked out, but a structural requirement for joint representation generation.
